\ifdefined\gkver\else\def\gkver{student}\fi
\documentclass[\gkver]{gaokaozhenti}
\begin{document}

\chapter{2008年上海理综}
%% paperType: 理综物理部分

\begin{enumerate}
\item
%% number: 4
%% paperName: 2008年上海理综第 4 题 3 分
%% typeId: 1
%% type: 单选题
%% chapter: 原子和原子核
%% point: 原子核式结构
%% method: 无
%% score: 3
%% degree: 850
%% duplicateId: 0
%% body:
二十世纪初，为了研究物质内部的结构，物理学家做了大量的实验，揭示了原子内部的结构，发现了电子、中子和质子，右图是\xzanswer{A}

\onepicture[w=4.2cm]{figs/04.png}

\fourchoices[answer=A]
{卢瑟福的 $\alpha$ 粒子散射实验装置}
{卢瑟福发现质子的实验装置}
{汤姆逊发现电子的实验装置}
{查德威克发现中子的实验装置}

%% answer: A
%% memo:
\memoanswer{
题目中所给装置是卢瑟福研究 $\alpha$ 粒子散射实验的装置，故选项 A 正确。
}

\item
%% number: 5
%% paperName: 2008年上海理综第 5 题 3 分
%% typeId: 1
%% type: 单选题
%% chapter: 气体性质
%% point: 盖吕萨克定律
%% method: 无
%% score: 3
%% degree: 750
%% duplicateId: 0
%% body:
温度计是生活、生产中常用的测温装置。右图为一个简易温度计，一根装有一小段有色水柱的细玻璃管穿过橡皮塞插入烧瓶内，封闭一定质量的气体。当外界温度发生变化时，水柱位置将上下变化。已知 A、D 间的测量范围为 $20\UdC\sim80\UdC$，A、D 间刻度均匀分布。由图可知，A、D 及有色水柱下端所示温度分别为\xzanswer{C}

\onepicture[w=1.8cm]{figs/05.png}

\fourchoices[answer=C]
{$20\UdC$、$80\UdC$、$64\UdC$}
{$20\UdC$、$80\UdC$、$68\UdC$}
{$80\UdC$、$20\UdC$、$32\UdC$}
{$80\UdC$、$20\UdC$、$34\UdC$}

%% answer: C
%% memo:
\memoanswer{
根据题意可知，温度越高，水柱上升的高度越高，A 点温度最高，D 点温度最低，故选项 A、B 错误。由于 A、D 间的刻度均匀分布，故水柱下端的温度为 $\frac{80-20}{15}\times3+20=32\UdC$，选项 C 正确。
}

\item
%% number: 6
%% paperName: 2008年上海理综第 6 题 3 分
%% typeId: 1
%% type: 单选题
%% chapter: 电磁感应
%% point: 电磁感应现象
%% method: 无
%% score: 3
%% degree: 750
%% duplicateId: 0
%% body:
老师做了一个物理小实验让学生观察：一轻质横杆两侧各固定一金属环，横杆可绕中心点自由转动，老师拿一条形磁铁插向其中一个小环，后又取出插向另一个小环，同学们看到的现象是\xzanswer{B}

\onepicture[w=3.2cm]{figs/06.png}

\fourchoices[answer=B]
{磁铁插向左环，横杆发生转动}
{磁铁插向右环，横杆发生转动}
{无论磁铁插向左环还是右环，横杆都不发生转动}
{无论磁铁插向左环还是右环，横杆都发生转动}

%% answer: B
%% memo:
\memoanswer{
左环没有闭合，在磁铁插入过程中，不产生感应电流，故横杆不发生转动。右环闭合，在磁铁插入过程中，产生感应电流，横杆将发生转动，故选项 B 正确。
}

\item
%% number: 7
%% paperName: 2008年上海理综第 7 题 3 分
%% typeId: 1
%% type: 单选题
%% chapter: 曲线运动
%% point: 天体运动
%% method: 无
%% score: 3
%% degree: 550
%% duplicateId: 0
%% body:
有同学这样探究太阳的密度：正午时分让太阳光垂直照射一个当中有小孔的黑纸板，接收屏上出现了一个小圆斑；测量小圆斑的直径和黑纸板到接收屏的距离，可大致推出太阳直径。他掌握的数据是：太阳光传到地球所需的时间、地球的公转周期、万有引力恒量；在最终得出太阳密度的过程中，他用到的物理规律是小孔成像规律和\xzanswer{C}

\fourchoices[answer=C]
{牛顿第二定律}
{万有引力定律}
{万有引力定律、牛顿第二定律}
{万有引力定律、牛顿第三定律}

%% answer: C
%% memo:
\memoanswer{
根据万有引力定律和牛顿第二定律 $G\frac{Mm}{r^{2}}=m\frac{v^{2}}{r}$ 可得太阳的质量，根据小孔成像规律和相似三角形的知识可得太阳的直径 $D$，故可求出太阳的密度，选项 C 正确。
}

\item
%% number: 8
%% paperName: 2008年上海理综第 8 题 3 分
%% typeId: 1
%% type: 单选题
%% chapter: 机械波
%% point: 波长频率波速
%% method: 无
%% score: 3
%% degree: 650
%% duplicateId: 0
%% body:
噪声会对人的心理、生理、生活与工作带来严重影响！通常用声强级 $L_{I}=10\lg\frac{I}{I_{0}}$（单位为 $\UdB$）来表示噪声的大小。式中，$I$ 为声强，单位是 $\UWmq$；$I_{0}=10^{-12}\UWmq$ 是人刚好能听到的声音强度。我国规定工作环境的噪声一般应低于 $85\UdB$，则以下最接近该标准的声强是\xzanswer{C}

\fourchoices[answer=C]
{$10^{-1}\UWmq$}
{$10^{-2}\UWmq$}
{$10^{-4}\UWmq$}
{$10^{-6}\UWmq$}

%% answer: C
%% memo:
\memoanswer{
根据题意有 $10\lg\frac{I}{10^{-12}}\leq85$，解得 $I\leq10^{-3.5}\UWmq\approx3.2\times10^{-4}\UWmq$，最接近 $10^{-4}\UWmq$，故选项 C 正确。
}

\item
%% number: 9
%% paperName: 2008年上海理综第 9 题 3 分
%% typeId: 1
%% type: 单选题
%% chapter: 光学
%% point: 光的电磁说
%% method: 无
%% score: 3
%% degree: 650
%% duplicateId: 0
%% body:
红外遥感卫星通过接收地面物体发出的红外辐射来探测地面物体的状况。地球大气中的水汽（\ce{H2O}）、二氧化碳（\ce{CO2}）能强烈吸收某些波长范围的红外辐射，即地面物体发出的某些波长的电磁波，只有一部分能够通过大气层被遥感卫星接收。右图为水和二氧化碳对某一波段不同波长电磁波的吸收情况，由图可知，在该波段红外遥感大致能够接收到的波长范围为\xzanswer{D}

\onepicture[w=6.5cm]{figs/09.png}

\fourchoices[answer=D]
{$2.5\sim3.5\Uum$}
{$4\sim4.5\Uum$}
{$5\sim7\Uum$}
{$8\sim13\Uum$}

%% answer: D
%% memo:
\memoanswer{
由图可知，$8\sim13\Uum$ 的波段被水和二氧化碳吸收的较少，能够接收到，故选项 D 正确。
}

\item
%% number: 11
%% paperName: 2008年上海理综第 11 题 3 分
%% typeId: 1
%% type: 单选题
%% chapter: 气体性质
%% point: 分子动理论
%% method: 无
%% score: 3
%% degree: 750
%% duplicateId: 0
%% body:
人们常用干冰营造云雾缭绕的舞台效果，这是因为干冰在室温下很容易直接变成气体。在此过程中体积可以增大很多倍，原因是\xzanswer{D}

\onepicture[w=8.5cm]{figs/11.png}

%% answer: D
%% memo:
\memoanswer{本题原卷及材料中未提供详解，待补充。}

\item
%% number: 42
%% paperName: 2008年上海理综第 42 题 2 分
%% typeId: 3
%% type: 填空题
%% chapter: 运动和力
%% point: 牛顿第二定律
%% method: 无
%% score: 2
%% degree: 950
%% duplicateId: 0
%% body:
汽车、火车、飞机在行驶中都会因某些因素而影响其速度，所以高速运行的交通工具的头部一般呈流线型，主要是为了减少 \tkanswer[1.5cm]{阻力}。

\onepicture[w=4.5cm, align=right]{figs/42.png}

%% answer:
\jdanswer{阻力}

%% memo:
\memoanswer{
根据流体知识可知，流线型可以减少阻力，故本题答案为阻力。
}

\item
%% number: 43
%% paperName: 2008年上海理综第 43 题 4 分
%% typeId: 3
%% type: 填空题
%% chapter: 运动和力
%% point: 牛顿第二定律
%% method: 无
%% score: 4
%% degree: 750
%% duplicateId: 0
%% body:
某汽车的部分参数如下表，请根据表中数据完成表的其他部分。

\begin{table}[htbp]
\centering
\begin{tblr}{|c|c|c|c|}
\hline
\SetCell[c=2]{c} 整车行驶质量 $1500\Ukg$ & & \SetCell[c=2]{c} 最大功率 $92\UkW$ & \\
\hline
\SetCell[r=2]{c} 加速性能 & \SetCell[c=2]{c} $0\sim108\Ukmh$（即 $30\Ums$）所需时间 & & 平均加速度 \\
\cline{2-4}
 & \SetCell[c=2]{c} $11\Us$ & & \tkanswer[1.8cm]{2.73} $\Umsq$ \\
\hline
\SetCell[r=2]{c} 制动性能 & \SetCell[c=2]{c} 车辆以 $36\Ukmh$（即 $10\Ums$）行驶时的制动距离 & & 制动过程中所受合外力 \\
\cline{2-4}
 & \SetCell[c=2]{c} $6.5\Um$ & & \tkanswer[1.8cm]{$1.15\times10^{4}$} $\UN$ \\
\hline
\end{tblr}
\end{table}

%% answer:
\jdanswer{2.73；$1.15\times10^{4}$}

%% memo:
\memoanswer{本题原卷及材料中未提供详解，待补充。}

\item
%% number: 44
%% paperName: 2008年上海理综第 44 题 3 分
%% typeId: 3
%% type: 填空题
%% chapter: 机械波
%% point: 多普勒效应
%% method: 无
%% score: 3
%% degree: 750
%% duplicateId: 0
%% body:
生活中经常用“呼啸而来”形容正在驶近的车辆，这是声波在传播过程中对接收者而言频率发生变化的表现。无线电波也具有这种效应。图中的测速雷达正在向一辆接近的车辆发出无线电波，并接收被车辆反射的无线电波。由于车辆的运动，接收的无线电波频率与发出时不同，利用频率差 $f_{\text{接收}}-f_{\text{发出}}$ 就能计算出车辆的速度。已知发出和接收的频率间关系为 $f_{\text{接收}}=\left(1+\frac{2v_{\text{车}}}{c}\right)f_{\text{发出}}$，式中 $C$ 是真空中的光速，若 $f_{\text{发出}}=2\times10^{9}\UHz$，$f_{\text{接收}}-f_{\text{发出}}=400\UHz$，可知被测车辆的速度大小为 \tkanswer[1.5cm]{30} $\Ums$。

\onepicture[w=4.2cm, align=right]{figs/44.png}

%% answer:
\jdanswer{30}

%% memo:
\memoanswer{
根据题意有 $\left(1+\frac{2v_{\text{车}}}{C}\right)f_{\text{发出}}-f_{\text{发出}}=400\UHz$，即 $\frac{2v_{\text{车}}}{C}f_{\text{发出}}=400\UHz$，解得 $v_{\text{车}}=\frac{400\times3\times10^{8}}{2\times2\times10^{9}}\Ums=30\Ums$。
}

\item
%% number: 45
%% paperName: 2008年上海理综第 45 题 3 分
%% typeId: 1
%% type: 单选题
%% chapter: 运动和力
%% point: 牛顿第二定律
%% method: 无
%% score: 3
%% degree: 650
%% duplicateId: 0
%% body:
有一个小实验：当向两片竖直放置的纸片的中间吹气时，会发现两个小纸片不但不分离，反而靠拢了。这一现象告诉我们，流体运动部分产生的压强要比它周围静止部分产生的压强小。也可以概括为流速大，压强小，流速小，压强大。飞机上天就是由于机翼上下空气流速不同造成的压力差所致。F1 赛车在尾部都装有尾翼板，其作用是\xzanswer{B}

\onepicture[w=4.2cm]{figs/45.png}

\fourchoices[answer=B]
{减小摩擦}
{抵消升力}
{减小阻力}
{升高重心}

%% answer: B
%% memo:
\memoanswer{
F1 赛车在尾部装有尾翼板，是为抵消升力，让赛车能稳定地在轨道上运动，故选项 B 正确。
}

\item
%% number: 46
%% paperName: 2008年上海理综第 46 题 4 分
%% typeId: 3
%% type: 填空题
%% chapter: 气体性质
%% point: 能源的开发与利用
%% method: 无
%% score: 4
%% degree: 750
%% duplicateId: 0
%% body:
电动车所需能量由它所携带的蓄电池供给。右图为某类电动车与汽油车性能的比较。通常车用电动机的质量是汽油机的 4 倍或 5 倍。为促进电动车的推广使用，在技术上主要应对电动车的 \tkanswer[1.5cm]{蓄电池}、\tkanswer[1.5cm]{电动机} 等部件加以改进。

给电动车蓄电池充电的能量实际上来自于发电站。一般发电站燃烧燃料所释放出来的能量仅 $30\%$ 转化为电能，在向用户输送及充电过程中又损失了 $20\%$，这意味着使用电动车时能量转化的总效率约为 \tkanswer[1.5cm]{16.8}$\%$。

\onepicture[w=7.5cm, align=right]{figs/46.png}

%% answer:
\jdanswer{蓄电池；电动机；16.8}

%% memo:
\memoanswer{
根据题意，使用电动车时能量转化的总效率约为 $30\%\times80\%\times70\%=16.8\%$。
}

\end{enumerate}

\end{document}
